\section*{Bab \ref{ch:b2:proba} --- Peluang pada Ruang Terbilang}

\begin{solution}{exo:b2:proba:1}
Berkat kesimetriannya: urutan penarikannya mengimbaskan urutan relatif
yang acak seragam pada bola $1$ dan $2$, jadi $\P(1 \text{ sebelum } 2) =
\frac12$. Secara formal: mempertukarkan kedudukan bola $1$ dan $2$
pada sebuah barisan penarikan merupakan bijeksi atas hasil
(yang berpeluang sama) yang menukar \omterm{def:b2:proba:space}{kejadiannya} dengan komplemennya. Adapun di antara bola
$1, \dots, k$: urutan relatif $k$ bolanya seragam
di antara $k!$ urutannya, dan bola $1$ menjadi yang pertama pada $(k-1)!$
di antaranya: jadi peluangnya $\frac{(k-1)!}{k!} = \frac1k$.
\end{solution}

\begin{solution}{exo:b2:proba:2}
Berlaku $\frac{1}{k(k+1)} = \frac1k - \frac1{k+1}$, jadi $\sum_{k\geq1} p_k$
berteleskop menjadi $1$: sehingga sebuah \omterm{def:b2:proba:space}{ukuran peluang}. Adapun hasil yang genap:
\[
\P(2\N^*) = \sum_{j=1}^{\infty}\frac{1}{2j(2j+1)}
= \sum_{j=1}^{\infty}\Bigl(\frac{1}{2j} - \frac{1}{2j+1}\Bigr)
= \frac12 - \frac13 + \frac14 - \frac15 + \cdots
\]
Ini adalah deret harmonik berselang-seling dengan suku pertamanya
dibuang dan tandanya dibalik: jadi karena $\ln 2 = 1 - \frac12 + \frac13 -
\frac14 + \cdots$ (\cref{ch:b2:series}),
\[
\P(2\N^*) = -\bigl(\ln 2 - 1\bigr) = 1 - \ln 2 \approx 0.307 .
\]
\end{solution}

\begin{solution}{exo:b2:proba:3}
Misalkan $S$ = sakit dan $+$ = uji positif. Maka Bayes
(\cref{thm:b2:proba:bayes}) dengan pemartisian $\{S, S^c\}$ memberi
\[
\P(S \mid +)
= \frac{0.99 \times 10^{-4}}
       {0.99 \times 10^{-4} + 0.01 \times 0.9999}
= \frac{0.000099}{0.000099 + 0.009999}
\approx 0.0098 ,
\]
yakni di bawah $1\%$. Jadi walaupun ujinya ``akurat 99\%'', sebuah hasil
positif menyisakan kamu sekitar $99\%$ mungkin sehat: sebab
positif palsu di antara mayoritas sehat yang sangat besar menenggelamkan
positif benar dari minoritas sakit yang sangat kecil. Jadi uji penapisan bagi
kondisi yang langka harus selalu dibaca lewat perhitungan laju-dasar
ini.
\end{solution}

\begin{solution}{exo:b2:proba:4}
\omterm{def:b2:funcseq:def}{Titik demi titik} pada $\Omega$: berlaku $\omega \in \bigcup A_i$ bila dan hanya bila suatu faktor
$1 - \mathbf{1}_{A_i}(\omega)$ lenyap, jadi
\[
\mathbf{1}_{\bigcup A_i}
= 1 - \prod_{i=1}^n\bigl(1 - \mathbf{1}_{A_i}\bigr)
= \sum_{\emptyset \neq J \subseteq \{1,\dots,n\}}
(-1)^{\abs J + 1}\prod_{i \in J}\mathbf{1}_{A_i} ,
\]
setelah hasil kalinya diuraikan lalu $1$-nya dipindahkan. Kini
$\prod_{i\in J}\mathbf{1}_{A_i} = \mathbf{1}_{\bigcap_{i \in J}
A_i}$, jadi menjumlahkannya terhadap bobot $\P(\{\omega\})$ ---
yang sah: sebab berhingga banyak suku yang terbatas, dan tiap keluarganya \omterm{def:b2:series:summable}{terjumlahkan} ---
mengubah tiap indikatornya menjadi peluang \omterm{def:b2:proba:space}{kejadiannya}, sehingga memberi
rumusnya.
\end{solution}

\begin{solution}{exo:b2:proba:5}
Misalkan $A_i$ = ``surat $i$ berada di amplop yang benar''. Untuk $J$
berukuran $k$ berlaku $\P\bigl(\bigcap_{i\in J}A_i\bigr) =
\frac{(n-k)!}{n!}$ (yakni tetapkan $k$ suratnya, lalu permutasikan sisanya). Jadi menurut
inklusi--eksklusi,
\[
\P\Bigl(\bigcup A_i\Bigr)
= \sum_{k=1}^n (-1)^{k+1}\binom nk \frac{(n-k)!}{n!}
= \sum_{k=1}^n \frac{(-1)^{k+1}}{k!} ,
\]
jadi
\[
\P(\text{tak ada kecocokan})
= 1 - \P\Bigl(\bigcup A_i\Bigr)
= \sum_{k=0}^{n}\frac{(-1)^k}{k!}
\xrightarrow[n\to\infty]{} e^{-1} \approx 0.368 .
\]
Adapun tepat satu kecocokan: sebuah permutasi dengan tepat satu titik tetap
ditentukan oleh pemilihan surat yang tetapnya (ada $n$ cara) dan sebuah
\emph{permutasi kacau} (yakni susunan tanpa kecocokan) atas $n - 1$ lainnya;
jadi dengan menulis $D_{n-1} = (n-1)!\sum_{k=0}^{n-1}\frac{(-1)^k}{k!}$ bagi
banyaknya permutasi kacau (yakni bagian pertamanya, yang diskalakan
$(n-1)!$),
\[
\P(\text{tepat satu kecocokan})
= \frac{n\,D_{n-1}}{n!}
= \frac{D_{n-1}}{(n-1)!}
= \sum_{k=0}^{n-1}\frac{(-1)^k}{k!}
\xrightarrow[n\to\infty]{} e^{-1} :
\]
pada limitnya, ``tak ada kecocokan'' dan ``tepat satu kecocokan'' sama
mungkinnya, masing-masing berpeluang $e^{-1}$.
\end{solution}

\begin{solution}{exo:b2:proba:6}
Syaratkanlah pada awalnya (yakni aturan rantai /
\cref{thm:b2:proba:bayes}):
\begin{itemize}
\item lemparan pertamanya A (berpeluang $1 - p$): permainannya bermula
ulang; jadi berlangsung lebih dari $n$ berarti berlangsung lebih dari $n - 1$ dari
sana: sehingga sumbangannya $(1-p)\,q_{n-1}$;
\item lemparan pertamanya GA (berpeluang $p(1-p)$): bermula ulang sesudah dua
lemparan: sumbangannya $p(1-p)\,q_{n-2}$;
\item lemparan pertamanya GG: permainannya telah berakhir (dalam $n$ lemparan, dengan $n
\geq 2$): jadi menyumbang $0$.
\end{itemize}
Karena itu $q_n = (1-p)q_{n-1} + p(1-p)q_{n-2}$. Adapun persamaan
karakteristiknya $r^2 = (1-p)r + p(1-p)$ berakar
\[
r_\pm = \frac{(1-p) \pm \sqrt{(1-p)^2 + 4p(1-p)}}{2},
\]
dengan $\abs{r_\pm} < 1$: sebab memang polinomial $\chi(r) = r^2 -
(1-p)r - p(1-p)$ memenuhi $\chi(1) = 1 - (1-p) - p(1-p) = p^2 >
0$ dan $\chi(-1) = 1 + (1-p) - p(1-p) > 0$, sedangkan $\chi(0) =
-p(1-p) < 0$: jadi satu akar di $\intoo{-1}{0}$ dan satu di $\intoo{0}{1}$.
Jadi $q_n = \alpha r_+^n + \beta r_-^n \to 0$. Adapun \omterm{def:b2:proba:space}{kejadian} ``permainannya
berlangsung lebih dari $n$'' turun ke ``permainannya tak pernah berakhir''; jadi kekontinuan
monoton (\cref{thm:b2:proba:continuity}) memberi
$\P(\text{tak pernah berakhir}) = \lim q_n = 0$: sehingga permainannya berakhir hampir
pasti.
\end{solution}

\begin{solution}{exo:b2:proba:7}
\emph{Bahwa $\P(R_n) = 1/n$:} di antara $n$ tarikan pertamanya, masing-masing $n$
kedudukan relatif tarikan terakhirnya sama mungkin (berkat keseragaman
urutan relatifnya), dan $R_n$ adalah \omterm{def:b2:proba:space}{kejadian} bahwa ia yang
terbesar: jadi peluangnya $1/n$.

\emph{Rekor tak berhingga banyaknya:} berlaku $\sum_n \P(R_n) = \sum 1/n =
\infty$ dan $R_n$ \omterm{def:b2:proba:independence}{saling bebas} (yang diterima saja), jadi
Borel--Cantelli~2 (\cref{thm:b2:proba:borelcantelli}) memberi
$\P(\limsup R_n) = 1$: jadi rekornya tak pernah berhenti, secara hampir pasti --- tetapi
ia menipis secara logaritmik.

\emph{Rekor berurutan:} berkat kebebasannya,
\[
\sum_n \P(R_n \cap R_{n+1})
= \sum_n \frac{1}{n(n+1)} < \infty ,
\]
jadi Borel--Cantelli~1 berlaku: sehingga hampir pasti, hanya berhingga kali
sebuah rekor langsung disusul rekor yang lain. Jadi kedua
paruh lemanya bekerja bersama: yakni rekor tak berhingga banyaknya, tetapi
(hampir pasti) akhirnya tak pernah dua berturut-turut.
\end{solution}

\begin{solution}{exo:b2:proba:8}
\emph{Bagian pertamanya:} berlaku $\sum \P(A_n) = \sum\frac{1}{n+1} = \infty$
dengan kebebasannya: jadi Borel--Cantelli~2 memberi $\P(\limsup A_n) = 1$.
Jadi tiap $A_n$ individualnya makin tak mungkin, namun hampir setiap
$\omega$ termasuk tak berhingga banyak di antaranya.

\emph{Contoh tandingan tanpa kebebasan:} ambillah $\Omega = \N^*$
dengan bobot $p_k = \frac{1}{k(k+1)}$ pada
\cref{exo:b2:proba:2}, dan $A_n = \{k \in \N^* : k \geq n\}$.
Maka
\[
\P(A_n) = \sum_{k \geq n}\Bigl(\frac1k - \frac1{k+1}\Bigr)
= \frac1n ,
\qquad
\sum_n \P(A_n) = \infty ,
\]
tetapi $A_n$ bersifat turun, jadi $\limsup_n A_n = \bigcap_n A_n =
\emptyset$: sehingga $\P(\limsup A_n) = 0$. Jadi kedivergenan $\sum\P(A_n)$
sendirian tak menjamin apa pun ketika \omterm{def:b2:proba:space}{kejadiannya} menumpuk pada bagian
ruangnya yang menciut --- dan kebebasanlah yang melarang
persekongkolan itu.
\end{solution}

\begin{solution}{exo:b2:proba:9}
Dengan $q = 1 - p$, gambar pertamanya jatuh pada pangkat $2j + 1$ dengan
peluang $q^{2j}p$, jadi
\[
\P(\text{pangkat ganjil}) = \sum_{j\geq0}q^{2j}p
= \frac{p}{1 - q^2} = \frac{1}{1 + q} .
\]
Untuk koin yang setimbang: berlaku $\frac1{1 + 1/2} = \frac23$. (Pemeriksaan kewarasannya:
pangkat ganjilnya sepatutnya lebih mungkin, karena pangkat $1$ datang lebih dulu ---
dan memang $\frac1{1+q} > \frac12$ selalu.)
\end{solution}

\begin{solution}{exo:b2:proba:10}
\omterm{def:b2:proba:space}{Kejadian} $B_N = \bigcap_{n=1}^N A_n^c$ turun ke
$\bigcap_nA_n^c$, dan berkat kebebasan komplemennya
berlaku $\P(B_N) = \prod_{n=1}^N(1 - p_n)$; jadi kekontinuan monoton
(\cref{thm:b2:proba:continuity}) memberi limit pada displainya.
Lalu dengan mengambil logaritmanya, $\prod(1 - p_n) > 0$ bila dan hanya bila $\sum-\ln(1 -
p_n) < \infty$. Bila $\sum p_n < \infty$ maka $p_n \to 0$ dan
$-\ln(1 - p_n) \sim p_n$: jadi deret lognya konvergen. Sedangkan bila
$\sum p_n = \infty$, maka $-\ln(1 - p_n) \geq p_n$ memaksa
kedivergenannya, jadi hasil kalinya $0$. Dan ini cocok dengan
Borel--Cantelli~2: sebab untuk $\sum p_n = \infty$, bukan hanya
$\P(\text{tak ada }A_n\text{ yang terjadi}) = 0$, melainkan hampir pasti
tak berhingga banyak $A_n$ terjadi.
\end{solution}

\begin{solution}{exo:b2:proba:11}
Katakanlah kotak $A$ yang pertama ditemukan kosong, dengan kotak
lainnya berisi $k$. Ini berarti: di antara $2n - k$
rogohan pertamanya, tepat $n$ menuju $A$ dan $n - k$ menuju $B$ (dalam suatu
urutan), lalu rogohan ke-$2n - k + 1$ menuju $A$ lagi
dan menemukannya kosong. Adapun rogohannya berupa pilihan setimbang yang \omterm{def:b2:proba:independence}{saling bebas},
jadi \omterm{def:b2:proba:space}{kejadian} ini berpeluang
$\binom{2n-k}{n}2^{-(2n-k)}\cdot\frac12$; lalu melipatduakannya (sebab kotak
kosongnya boleh yang mana saja) memberi
\[
\P(\text{kotak lainnya berisi }k) = \binom{2n-k}{n}\,2^{-(2n-k)} .
\]
Untuk $n = 1$: $k = 1$ memberi $\binom11 2^{-1} = \frac12$ dan $k
= 0$ memberi $\binom21 2^{-2} = \frac12$: jadi totalnya $1$, sebagaimana
mestinya.
\end{solution}

\begin{solution}{exo:b2:proba:12}
(a) Bila $\P(\{n\}) = c$ untuk setiap $n$, maka keaditifan-$\sigma$
memaksa $1 = \sum_nc$: yang mustahil, entah $c = 0$ (sebab jumlahnya $0$)
atau $c > 0$ (sebab jumlahnya tak hingga). Jadi tak ada peluang seragam pada
$\N$.

(b) Bila $A \cap B = \emptyset$ dan $d(A)$ serta $d(B)$ ada, maka
$\abs{(A \sqcup B)\cap\intint1n} = \abs{A\cap\intint1n} +
\abs{B\cap\intint1n}$, jadi $d(A \sqcup B) = d(A) + d(B)$:
yakni keaditifan berhingga pada pasangan semacam itu. Adapun tiap singletonnya punya fungsi
pencacah yang akhirnya tetap, jadi berkerapatan $0$, sedangkan
$d(\N^*) = 1$. Seandainya $d$ bersifat aditif-$\sigma$, maka $\N^* =
\bigsqcup_k\{k\}$ akan memberi $1 = \sum_k 0 = 0$: jadi kerapatannya
aditif berhingga tetapi tak aditif-$\sigma$ --- sehingga aksiomanya punya
isi.

(c) Misalkan $A = \bigcup_{k\geq0}\intint{4^k}{2\cdot4^k - 1}$
(yakni blok dari $4^k$ sampai $2\cdot4^k - 1$). Di $n = 2\cdot4^K -
1$ cacahnya $\sum_{k\leq K}4^k \sim \frac43 4^K$, sehingga memberi
nisbah $\to \frac23$; sedangkan di $n = 4^{K+1} - 1$ cacahnya
tak berubah, sehingga memberi nisbah $\to \frac13$. Jadi nisbahnya berayun
antara limit $\frac13$ dan $\frac23$: sehingga tak berkerapatan.
\end{solution}

\begin{solution}{pb:b2:proba:1}
\textbf{1.} Sebuah lintasan berpanjang $n$ ditentukan oleh himpunan
langkah naiknya; dan berakhir di $k$ berarti ada $u$ langkah naik dan $n - u$
langkah turun dengan $u - (n - u) = k$, yakni $u = \frac{n+k}2$:
yang mungkin bila dan hanya bila $n + k$ genap dan $\abs k \leq n$, dalam
$\binom{n}{(n+k)/2}$ cara. Adapun tiap lintasan tertentunya berupa satu titik
ukuran hasil kali setimbang pada $n$ lemparan: jadi berpeluang $2^{-n}$.
Karena itu $\P(S_n = k) = N_n(k)2^{-n}$.

\textbf{2.} Peubah $S_n$ berparitas $n$, jadi $S_{2n+1} \neq
0$; dan $u_n = N_{2n}(0)4^{-n} = \binom{2n}n4^{-n}$. Adapun nilainya:
$u_1 = \frac12$, $u_2 = \frac6{16} = \frac38$, dan $u_3 =
\frac{20}{64} = \frac5{16}$.

\textbf{3.} Berlaku $\dfrac{u_n}{u_{n-1}} =
\dfrac{\binom{2n}n}{4\binom{2n-2}{n-1}} =
\dfrac{(2n)(2n-1)}{4n^2} = \dfrac{2n-1}{2n} < 1$: jadi turun.
Lalu menurut \cref{ex:b2:comparison:centralbinomial}, $\binom{2n}n \sim
\frac{4^n}{\sqrt{\pi n}}$, jadi $u_n \sim \frac1{\sqrt{\pi n}}
\to 0$, dan $\sum u_n$ divergen lewat pembandingan dengan $\sum
n^{-1/2}$.

\textbf{4.} Diberikan sebuah lintasan dari $1$ ke $k$ yang menyentuh $0$,
pantulkanlah ruas awalnya (sampai kunjungan \emph{pertamanya} ke
$0$) terhadap sumbu mendatarnya: hasilnya sebuah lintasan dari
$-1$ ke $k$, dan operasinya berupa involusi --- sebab setiap
lintasan dari $-1$ ke $k \geq 1$ mesti menyeberangi $0$, dan memantulkan
ruas awalnya kembali memulihkan aslinya. Karena itu lintasan
yang menyentuhnya berjumlah $N_{n-1}(k + 1)$ (sebab dari $-1$ ke $k$
perpindahannya $k + 1$). Adapun lintasan dari $0$ ke $k$ yang tinggal $>
0$ sesudah waktu $0$ bermula dengan langkah naik lalu berjalan dari
$1$ ke $k$ dalam $n - 1$ langkah tanpa menyentuh $0$: jadi ada
$N_{n-1}(k-1) - N_{n-1}(k+1)$ di antaranya.

\textbf{5.} Dengan $m = \frac{n+k}2$, dan memakai
$\binom{n-1}{m-1} = \frac mn\binom nm$ serta $\binom{n-1}{m} =
\frac{n-m}n\binom nm$:
\[
\frac{N_{n-1}(k-1) - N_{n-1}(k+1)}{N_n(k)}
= \frac{\binom{n-1}{m-1} - \binom{n-1}{m}}{\binom nm}
= \frac{m - (n - m)}{n} = \frac kn .
\]
Untuk $n = 3$ dan $k = 1$: ada $N_3(1) = 3$ lintasan (yakni $++-$, $+-+$,
$-++$), dan hanya $++-$ yang tinggal positif (sebab $+-+$ pulang ke
$0$ pada waktu $2$): jadi satu dari tiga, dan $\frac kn = \frac13$.

\textbf{6.} Berkat kesimetriannya peluangnya adalah $2\P(S_i > 0\ \forall
i \leq 2n)$. Lalu menjumlahkannya atas titik ujungnya $2k$ dan memakai
pertanyaan 4 (dengan $n$ diganti $2n$):
\[
\P(S_i > 0\ \forall i) = 2^{-2n}\sum_{k\geq1}
\bigl(N_{2n-1}(2k-1) - N_{2n-1}(2k+1)\bigr)
= 2^{-2n}\,N_{2n-1}(1),
\]
yakni sebuah jumlah yang berteleskop. Kini $N_{2n-1}(1) = \binom{2n-1}{n}$ dan
$2\binom{2n-1}n = \binom{2n}n$ (menurut Pascal), jadi peluang pada
displainya adalah $2\cdot2^{-2n}\binom{2n-1}n =
\binom{2n}n4^{-n} = u_n$.

\textbf{7.} \omterm{def:b2:proba:space}{Kejadian} $D_n = \{S_i \neq 0,\ i \leq 2n\}$
turun, dengan irisannya ``tak pernah pulang''; jadi menurut kekontinuan
monoton dan pertanyaan 6, $\P(\text{tak ada kepulangan}) = \lim u_n =
0$: sehingga jalannya pulang hampir pasti. Lebih jauh $f_n = \P(D_{n-1})
- \P(D_n) = u_{n-1} - u_n$, dan menurut pertanyaan 3
\[
u_{n-1} - u_n = u_n\Bigl(\frac{2n}{2n-1} - 1\Bigr) =
\frac{u_n}{2n-1};
\qquad
\sum_{n\geq1}f_n = u_0 - \lim u_n = 1 .
\]

\textbf{8.} Berlaku $2n\,f_n = \frac{2n}{2n-1}u_n \geq u_n$, dan
$\sum u_n = \infty$ (pertanyaan 3): jadi deret $\sum 2nf_n$
divergen. Sehingga kepulangan pertamanya pasti tetapi tak punya waktu tunggu
rerata yang berhingga --- jadi jalannya bersifat \emph{rekuren nol}, dalam
kosakata yang kelak disediakan \cref{ch:b2:randomvar}.

\textbf{9.} \omterm{def:b2:proba:space}{Kejadian} ``sekurangnya $k$ kepulangan'' merupakan
gabungan \omterm{def:b2:structures:countable}{terbilang} yang lepas, atas $0 < n_1 < \dots < n_k$, bagi
\omterm{def:b2:proba:space}{kejadian} ``$k$ kepulangan pertamanya terjadi tepat pada waktu
$2n_1, \dots, 2n_k$''. Adapun \omterm{def:b2:proba:space}{kejadian} semacam itu adalah irisan
$k$ \omterm{def:b2:proba:space}{kejadian} yang bergantung pada blok lemparan yang saling lepas
$\intint1{2n_1}$, $\intint{2n_1+1}{2n_2}$, \dots, dan tiap bloknya
menuntut sebuah jalan baru membuat kepulangan pertamanya sesudah
tepat sebanyak langkah yang dijatahkan; jadi berkat kebebasan
bloknya peluangnya adalah $f_{n_1}f_{n_2-n_1}\cdots
f_{n_k-n_{k-1}}$. Lalu setelah dijumlahkan lewat paket (\cref{ch:b2:series},
sebab semua sukunya taknegatif):
\[
\P(\text{sekurangnya }k\text{ kepulangan})
= \Bigl(\sum_{n\geq1}f_n\Bigr)^{\!k} = 1^k = 1 .
\]
\omterm{def:b2:proba:space}{Kejadiannya} turun dalam $k$, jadi kekontinuan monoton memberi
$\P(\text{pulang tak hingga kali}) = 1$: itulah kerekurenannya.

\textbf{10.} Menurut pertanyaan 9, jalannya membuat tak berhingga banyak
ekskursi menjauh dari $0$. Adapun langkah pertama tiap ekskursinya berupa
koin baru yang bebas dari segala sebelumnya: jadi
peluang bahwa $m$ ekskursi pertamanya semuanya bermula ke bawah
adalah $2^{-m}$. Lalu untuk mencapai $1$ jalannya hanya menuntut satu ekskursi
yang bermula ke atas (sebab dari $<0$ ia mesti melewati $0$ sebelum
mencapai $1$, karena langkahnya $\pm1$), jadi $\P(\text{tak pernah mengenai }1)
\leq 2^{-m}$ untuk setiap $m$: sehingga jalannya mengenai $1$ hampir pasti.
Lalu dengan menguraikannya atas waktu pengenaannya (yang hampir pasti berhingga),
jalan yang dimulai ulang di sana merupakan jalan baru yang bermula di $1$: jadi
secara induktif ia mengenai setiap $k \geq 1$ hampir pasti, dan berkat
kesimetriannya setiap $k \leq -1$. Akhirnya, dengan memulai ulang di kunjungan
pertamanya ke $k$, pertanyaan 9 berlaku bagi jalan barunya: jadi setiap
tapaknya dikunjungi tak berhingga sering, dan itu secara hampir pasti.

\textbf{11.} \omterm{def:b2:proba:space}{Kejadian} $A_n = \{S_{2n} = 0\}$ jauh dari
\omterm{def:b2:proba:independence}{saling bebas} (sebab berada di $0$ pada waktu $2n$ membuat berada di $0$ pada
waktu $2n + 2$ jauh lebih mungkin daripada $u_{n+1}$), jadi
Borel--Cantelli~2 tak tersedia, dan memang seluruh kerja
Bagian II adalah untuk menggantikannya. Sedangkan arah lainnya tak menuntut
kebebasan: sebab \emph{bila} $\sum\P(A_n)$ konvergen, maka
Borel--Cantelli~1 menghasilkan berhingga banyak kepulangan hampir pasti.
Dan implikasi itulah mesin setiap bukti ketransienan
di bawah.

\textbf{12.} Sebuah kepulangan pada waktu $2n$ menuntut $n$ langkah naik dan $n$
langkah turun: jadi $\P(S_{2n} = 0) = \binom{2n}np^nq^n =
u_n(4pq)^n$, dan $4pq = 1 - (p - q)^2 < 1$ untuk $p \neq
\frac12$. Lalu karena $u_n \leq 1$, deret $\sum\P(S_{2n} = 0)$
didominasi deret geometri $\sum(4pq)^n$: jadi konvergen. Sehingga menurut
Borel--Cantelli~1, $\P(S_{2n} = 0 \text{ tak berhingga sering}) =
0$: jadi berhingga banyak kepulangan, secara hampir pasti.

\textbf{13.} Untuk $n + k$ yang genap berlaku $\P(S_n = k) =
\binom{n}{\frac{n+k}2}p^{\frac{n+k}2}q^{\frac{n-k}2}$; adapun
koefisien binomialnya paling banyak sebesar yang pusatnya, dan
$p^{\frac{n+k}2}q^{\frac{n-k}2} = (pq)^{n/2}(p/q)^{k/2}$,
sehingga memberi batas yang dinyatakan $\leq
2^n(pq)^{n/2}(p/q)^{k/2} = (4pq)^{n/2}(p/q)^{k/2}$, yang \omterm{def:b2:series:summable}{terjumlahkan}
dalam $n$ karena $\sqrt{4pq} < 1$. Jadi Borel--Cantelli~1: tapak $k$
dikunjungi berhingga sering hampir pasti; lalu gabungan atas $k \in
\Z$ bagi \omterm{def:b2:proba:space}{kejadian} nol kekecualiannya tetap nol (berkat kesubaditifan
\omterm{def:b2:structures:countable}{terbilangnya}). Jadi hampir pasti setiap tapaknya dikunjungi berhingga
sering, sehingga barisan bulat $(S_n)$ meninggalkan setiap jendela terbatasnya
untuk selamanya: yakni $\abs{S_n} \to \infty$.

\textbf{14.} Berlaku $\P(S_n \neq 0,\ 1 \leq n \leq 200) = u_{100} =
\binom{200}{100}4^{-100} \approx \frac1{\sqrt{100\pi}}
\approx 0.056$: jadi lebih dari satu peluang dari dua puluh bahwa $200$
lemparan setimbang tak pernah seri. Adapun peluruhan $1/\sqrt{\pi n}$-nya
menyiksa saking lambatnya: sebab kepastian seri (pertanyaan 7)
serasi dengan rentang tanpa seri yang sangat \omterm{def:b2:curves:length}{panjang} --- yakni cita
rasa pertama gejala arcsinus pada Bagian IV.

\textbf{15.} Syaratkanlah pada langkah pertamanya. Bila $X_1 = +1$ maka
$T_1 = 1$, dan $f_1 = \frac12$ cocok. Sedangkan bila $X_1 = -1$,
jalannya harus mendaki dari $-1$ ke $1$; jadi menurut penguraian bloknya,
pulang ke $0$ untuk pertama kalinya pada waktu $2n$ terpecah sebagai:
satu langkah turun, lalu sebuah jalan baru yang bermula di $-1$ pertama kali
mencapai $0$ --- yang setara dengan sebuah jalan baru yang pertama kali mencapai
$+1$ --- dalam $2n - 1$ langkah, atau \omterm{def:b2:proba:space}{kejadian} setangkupnya ke atas.
Kedua tandanya menyumbang sama:
\[
f_n = 2\cdot\tfrac12\,\P(T_1 = 2n - 1) = \P(T_1 = 2n-1) .
\]
Karena itu $\P(T_1 < \infty) = \sum f_n = 1$, sedangkan $\sum_n(2n -
1)f_n = \sum_n u_n = \infty$ menurut pertanyaan 7: jadi jalannya mencapai
$1$ hampir pasti, dalam waktu rerata yang tak hingga.

\textbf{16.} Partisikanlah $\{M_n \geq k\}$ menurut nilai akhirnya
$S_n = m$. Untuk $m \geq k$ syarat $M_n \geq k$ berlaku
otomatis. Sedangkan untuk $m < k$, pantulkanlah lintasannya sesudah
kunjungan \emph{pertamanya} ke aras $k$: ini merupakan bijeksi antara
$\{M_n \geq k, S_n = m\}$ dan $\{S_n = 2k - m\}$ (sebab setiap lintasan
yang berakhir di $2k - m > k$ mengunjungi $k$; dan memantulkannya kembali adalah
baliknya). Karena itu
\[
\P(M_n \geq k)
= \sum_{m > k}\P(S_n = m) + \P(S_n = k)
+ \sum_{m < k}\P(S_n = 2k - m)
= 2\P(S_n > k) + \P(S_n = k).
\]

\textbf{17.} Pada waktu genap $2n$ dengan $k = 1$: berlaku $\P(S_{2n} = 1)
= 0$ dan $\P(S_{2n} > 1) = \P(S_{2n} \geq 2)$, jadi
\[
\P(M_{2n} \geq 1) = 2\P(S_{2n} \geq 2)
= \P(S_{2n} \geq 2) + \P(S_{2n} \leq -2)
= 1 - u_n .
\]
Jadi $\P(S_i \leq 0\ \forall i \leq 2n) = u_n$: sehingga jalannya
tak pernah unggul pada $2n$ langkah pertamanya sesering ia
tak pernah seri (pertanyaan 6) --- yakni dua \omterm{def:b2:proba:space}{kejadian} yang cukup berbeda,
yang diusung $u_n$ yang sama.

\textbf{18.} Berlaku $\{L_{2n} = 2k\} = \{S_{2k} = 0\} \cap
\{\text{jalan atas lemparan } 2k+1, \dots, 2n \text{ tak
bernol}\}$. Adapun kedua \omterm{def:b2:proba:space}{kejadiannya} bergantung pada blok lemparan yang saling lepas,
jadi keduanya \omterm{def:b2:proba:independence}{saling bebas}; yang pertama berpeluang $u_k$, dan yang
kedua $u_{n-k}$ menurut pertanyaan 6 yang diterapkan pada jalan baru
berlangkah $(2n-2k)$. Karena itu $\P(L_{2n} = 2k) = u_ku_{n-k}$.
Lalu karena $L_{2n}$ mengambil persis nilai $0, 2, \dots, 2n$,
peluang itu berjumlah $1$:
yakni $\sum_{k=0}^nu_ku_{n-k} = 1$, sebuah kesamaan binomial yang diserahkan
oleh sebuah pemartisian peluang.

\textbf{19.} Kesetangkupannya seketika: sebab $u_ku_{n-k} =
u_{n-k}u_k$. Lalu karena $u_j$ turun dalam $j$, hasil kali
$u_ku_{n-k}$ terkecil bagi $k$ yang di tengah dan terbesar di
ekstremnya $k \in \{0, n\}$, tempat ia sama dengan $u_n$;
secara kuantitatif $u_ku_{n-k} \approx \frac1{\pi\sqrt{k(n-k)}}$
di tengahnya, terhadap $u_n \approx \frac1{\sqrt{\pi n}}$ di
tepinya. Untuk $n = 5$: $\P(L_{10} = 0) = \P(L_{10} = 10) =
u_5 = \frac{63}{256} \approx 0.246$, sedangkan $\P(L_{10} = 4) =
u_2u_3 = \frac38\cdot\frac5{16} = \frac{15}{128} \approx
0.117$. Jadi pada permainan setimbang yang \omterm{def:b2:curves:length}{panjang}, penyamaan terakhirnya paling
mungkin di dekat awal sekali atau akhir sekali: sehingga satu pemain
biasanya unggul untuk rentang yang sangat \omterm{def:b2:curves:length}{panjang}, tanpa keberatsebelahan pada
koinnya.

\textbf{20.} Gambarannya: pada waktu $n$ jalannya hidup pada skala
$\sqrt n$ (yakni rentang binomial pertanyaan 3 --- sebab $u_n \sim
1/\sqrt{\pi n}$ adalah tinggi puncak pusatnya); lalu ia
pulang ke $0$ tak berhingga sering dengan peluang $1$ (Bagian
II), namun waktu tunggu di antara kepulangannya bererata divergen
(pertanyaan 8), dan itulah sebabnya satu ekskursi dapat menempati
pecahan positif sembarang cakrawala; sepadan dengan itu seri terakhir
sebuah permainan berlangkah $2n$ tersebar dengan nilai ekstremnya
paling mungkin (pertanyaan 18--19), dan tak pernah unggul berpeluang
$u_n$ yang meluruh lambat sama seperti tak pernah seri
(pertanyaan 17). Jadi kepastian pada limitnya, keawetan pada setiap
cakrawala berhingganya: itulah jalan yang setimbang.

\textbf{21.} Partisikanlah $\{S_{2n} = 0\}$ (dengan $n \geq 1$) menurut
waktu kepulangan pertamanya $2k$ dengan $1 \leq k \leq n$: maka blok pertama berisi
$2k$ lemparan mewujudkan sebuah kepulangan pertama, sedangkan $2n - 2k$
lemparan sisanya mewujudkan kepulangan sebuah jalan baru, dan bloknya
\omterm{def:b2:proba:independence}{saling bebas}: jadi $u_n = \sum_{k=1}^nf_ku_{n-k}$. Adapun kedua deretnya
$U(x) = \sum u_nx^n$ dan $F(x) = \sum f_nx^n$ berjari-jari $\geq
1$ (sebab koefisiennya di $\intcc01$), dan \omterm{thm:b2:series:fubini}{hasil kali Cauchynya}
(\cref{ch:b2:powerseries}) memberi, untuk $0 \leq x < 1$,
\[
U(x) - 1 = \sum_{n\geq1}\Bigl(\sum_{k=1}^n
f_ku_{n-k}\Bigr)x^n = F(x)\,U(x),
\qquad\text{yakni}\qquad
U(x)\bigl(1 - F(x)\bigr) = 1 .
\]

\textbf{22.} Ketika $x \uparrow 1$, fungsi $U(x)$ dan $F(x)$ naik
(sebab koefisiennya taknegatif); dan tiap jumlah parsial $\sum_{n\leq
N}u_n$ merupakan limit $\sum_{n\leq N}u_nx^n \leq U(x)$, jadi
$U(x) \uparrow \sum u_n \in \intoc0{+\infty}$, dan demikian pula
$F(x) \uparrow f = \sum f_n$. Bila $\sum u_n = \infty$: maka $1 -
F(x) = 1/U(x) \to 0$, jadi $f = 1$. Sedangkan bila $\sum u_n = S < \infty$:
maka $1 - f = 1/S > 0$, jadi $f < 1$. Adapun pemeriksaannya: jalan setimbangnya, $\sum u_n =
\infty$ dan $f = 1$ (pertanyaan 3, 7); sedangkan jalan berat sebelahnya, $\sum
u_n(4pq)^n < \infty$ dan sepadan dengan itu $f = 1 -
1/\sum_{n\geq0}u_n(4pq)^n < 1$, yang konsisten dengan
keberhinggaan hampir pasti banyaknya kepulangan (pertanyaan
12).

\textbf{23.} Untuk keempat langkah $(\pm1, 0), (0, \pm1)$ pada
jalan $\Z^2$, pertambahan $U = X + Y$ dan $V = X - Y$
adalah: $(+,+)$ bagi $(1,0)$, $(+,-)$ bagi $(0,1)$, $(-,+)$ bagi
$(0,-1)$, dan $(-,-)$ bagi $(-1,0)$ --- dengan tiap pasangan tandanya
berpeluang $\frac14 = \frac12\cdot\frac12$: jadi kedua
jalan koordinatnya $(U_n)$ dan $(V_n)$ merupakan jalan setimbang yang \omterm{def:b2:proba:independence}{saling bebas}
pada $\Z$. Lalu karena $S^{(2)}_{2n} = (0,0)$ bila dan hanya bila $U_{2n} = 0$
dan $V_{2n} = 0$,
\[
\P\bigl(S^{(2)}_{2n} = (0,0)\bigr) = u_n^2 \sim \frac1{\pi
n}, \qquad \sum_nu_n^2 = \infty .
\]
Adapun kesamaan pembaruan pertanyaan 21 dan dikotomi
pertanyaan 22 tak memakai apa pun yang satu dimensi (melainkan hanya
penguraian atas kepulangan pertamanya beserta kebebasan blok yang
saling lepas), jadi $\sum u_n^{(2)} = \infty$ memberi $f^{(2)} =
1$, lalu hujah pertanyaan 9 meningkatkannya: sehingga jalan pada
$\Z^2$ pulang ke titik asalnya tak berhingga sering hampir pasti.

\textbf{24.} Dengan batas yang diterima $\P(S^{(3)}_{2n} = 0)
\leq Cn^{-3/2}$, deretnya konvergen, jadi Borel--Cantelli~1
memberi berhingga banyak kepulangan hampir pasti: sehingga jalan pada
$\Z^3$ bersifat transien (dan batas yang sama dengan eksponen $-d/2$
menangani setiap $d \geq 3$). Jadi semuanya: yakni \emph{teorema
Pólya} --- bahwa \omterm{pb:b2:proba:1}{jalan acak sederhana} bersifat rekuren pada $\Z$ dan
$\Z^2$, serta transien pada $\Z^d$ untuk $d \geq 3$. Seorang pemabuk menemukan
jalan pulangnya; sedangkan seekor burung yang mabuk boleh jadi tidak.

\textbf{25.} Pencacahan lintasan dan pemantulan menghasilkan hukum
\omterm{def:b2:multint:exact}{eksaknya} ($u_n$, teorema pemungutan suara, $f_n$, maksimumnya, nol
terakhirnya); kekontinuan monoton mengubah setiap pernyataan limit
(yakni ``pulang sekurangnya sekali'' dan ``tak berhingga sering'')
menjadi limit peluang bercakrawala berhingga;
kebebasan blok yang saling lepas menggerakkan penguraian
pembaruannya (pertanyaan 9, 18, 21) --- sebab ia rangka
\omterm{def:b2:structures:countable}{terbilang} bagi sifat Markovnya; Borel--Cantelli~1 menjadi
senjata ketransienannya (pertanyaan 12--13, 24), yang tak menuntut
kebebasan; dan kesamaan pembaruannya menata semuanya menjadi
dikotomi $\sum u_n = \infty \iff$ kerekurenan. Adapun satu
masukan analitisnya adalah taksiran lokal $u_n \sim 1/\sqrt{\pi
n}$: sebab kuadratnya $1/(\pi n)$ masih divergen (jadi dimensi $2$
rekuren), sedangkan $n^{-3/2}$ konvergen (jadi dimensi $3$
transien) --- sehingga teorema Pólya, pada akhirnya, adalah pernyataan
tentang kedivergenan $\sum n^{-d/2}$.

\end{solution}
